% =====================================================================
%  Supplementary Material for the RECOMB submission.
%
%  RECOMB takes the appendix as a SEPARATE PDF with no page limit, so this is
%  built on its own and is not \input by recomb.tex:
%      latexmk -pdf appendix.tex
%
%  Tables come from paper/appendix_tables.tex, which is generated by
%  analysis/make_appendix_tables.py straight from the result files. Rerun that
%  script after any experiment changes; do not edit the tables by hand.
% =====================================================================

\documentclass[10pt]{article}
\usepackage[letterpaper,margin=1in]{geometry}
\def\RECOMB{1}

\usepackage[utf8]{inputenc}
\usepackage[T1]{fontenc}
\usepackage{lmodern}      % scalable font: microtype expansion needs one
\usepackage{microtype}
\usepackage{graphicx}
\usepackage{booktabs}
\usepackage{amssymb}
\usepackage[round]{natbib}
\usepackage[skip=4pt,font=small,labelfont=bf]{caption}
\usepackage{enumitem}
\usepackage{titling}
\usepackage{titlesec}
\usepackage{xcolor}       % hyperref's own color pkg has no blend syntax
\usepackage[colorlinks=true,allcolors=blue!60!black]{hyperref}

% compact lists -- a 4-page budget does not pay for list padding
\setlist[itemize]{topsep=2pt,itemsep=1pt,parsep=0pt,leftmargin=1.2em}
\setlength{\parskip}{0pt}
\graphicspath{{figures/}}

% [VENUE] the venue's style sets its own title block -- delete this group when
% you swap one in. article's default wastes ~1.5in at the top of page 1.
% \preauthor/\postauthor are overridden to keep the byline to one compact line;
% article's defaults cost about 8 more lines and push the body onto page 5.
\ifdefined\RECOMB
  \setlength{\droptitle}{-2em}
\else
  \setlength{\droptitle}{-6em}   % 4-page build needs the space
\fi
\pretitle{\begin{center}\LARGE}
\posttitle{\par\end{center}\vskip 0.4em}
\preauthor{\begin{center}\normalsize}
\postauthor{\par\end{center}}
\predate{}\postdate{}
\setlength{\abovedisplayskip}{4pt}
\setlength{\belowdisplayskip}{4pt}
\titlespacing*{\section}{0pt}{1.6ex plus .2ex}{0.9ex}
\titlespacing*{\paragraph}{0pt}{1.0ex plus .2ex}{0.6em}

% float spacing: article's defaults leave a wide band below a float and a narrow
% one above it, which reads as the caption belonging to the text underneath
\setlength{\textfloatsep}{14pt plus 2pt minus 2pt}   % float <-> text, top/bottom of page
\setlength{\floatsep}{12pt plus 2pt minus 2pt}       % between stacked floats
\setlength{\intextsep}{12pt plus 2pt minus 2pt}      % float placed mid-text

% float placement: six floats cluster in sections 2-4, and article's defaults
% (2 per top, 3 per page, 70% of the page) make them queue up and drift forwards
% into later sections, landing inside other sections' lists.
\setcounter{topnumber}{3}
\setcounter{bottomnumber}{2}
\setcounter{totalnumber}{4}
\renewcommand{\topfraction}{0.85}
\renewcommand{\bottomfraction}{0.6}
\renewcommand{\textfraction}{0.1}
\renewcommand{\floatpagefraction}{0.7}

% float-only pages: article's defaults are stretchy glue, which spreads two
% figures to opposite ends of the page. Pack them to the top instead.
\makeatletter
\setlength{\@fptop}{0pt}
\setlength{\@fpsep}{14pt}
\setlength{\@fpbot}{0pt plus 1fil}
\makeatother
\setlength{\abovecaptionskip}{6pt}
\setlength{\belowcaptionskip}{0pt}

\newcommand{\refalt}{\texttt{refalt}}
\newcommand{\refcopy}{\texttt{ref\_copy}}
\newcommand{\diffonly}{\texttt{diff\_only}}

\title{Long-context DNA models do not use the allele:\\
an audit of DNALongBench's eQTL task}

% [VENUE] anonymise for submission if the venue is double-blind: comment out the
% \author line below and restore \preauthor{}\postauthor{} above.
\ifdefined\RECOMB
  \author{Karandeep Nanda\\\normalsize University of Colorado Boulder}
\else
  \author{Karandeep Nanda, \normalsize University of Colorado Boulder}
\fi
\date{}



\title{Supplementary Material\\
\large Long-context DNA models do not use the allele:\\
an audit of DNALongBench's eQTL task}

\begin{document}
\maketitle

\noindent
This supplement reports the per-tissue, per-fold and per-seed numbers behind the
main text, the construction of eQTLP-v2 and the clinical set, and the released-code
defects with file and line references. Every table is generated from the result
files by \texttt{analysis/make\_appendix\_tables.py}.

\tableofcontents

\section{The published protocol, per tissue}
\label{app:a}

Table~\ref{tab:a1} gives the per-tissue view of the published eQTLP split that
\S2 of the main text summarises.

Three things to read off it. First, \textbf{distance alone} scores 0.688--0.828
per tissue against published foundation-model numbers of 0.48--0.58: the
zero-training baseline wins in every tissue. Second, the \textbf{test sets are
tiny} --- 11--34 positives --- so the 95\% CIs on the distance baseline are
0.10--0.19 wide, and no published FM number is separable from 0.5 at this sample
size. Third, the \textbf{split is not stratified and not disjoint}: positive
rates fall from 9--21\% in train to 3--5\% in test, every test gene already
appears in training, and 82--92\% of test pairs have a training variant within
100\,kb.

We score distance from the released coordinates, $|\mathrm{variant} -
\mathrm{TSS}|$, because that is what the padded input encodes: each input is
right-padded with N to 450\,kb, so the count of N tokens is a direct readout of
the span. The released \texttt{distance\_to\_tss} column disagrees with the
coordinates by more than 1\,kb for 16\% of negatives and no positives; scoring by
that column instead gives a mean of 0.663, still above every FM.

% generated by analysis/make_appendix_tables.py -- do not edit by hand
\begin{table}[tb]
\centering
\caption{Published eQTLP test split, per tissue. \textbf{dist} is our zero-training distance baseline scored from the released coordinates, with a 95\% bootstrap CI; \textbf{expert}, \textbf{hyena} and \textbf{cad-ph} are the published numbers. Every test gene appears in training.}
\label{tab:a1}
\footnotesize
\setlength{\tabcolsep}{3pt}
\begin{tabular}{lccccccccccc}
\toprule
Tissue & pos & neg & $p$ train & $p$ test & dist & 95\% CI & expert & hyena & cad-ph & gene & locus \\
\midrule
Adipose\_Subcutaneous & 16 & 316 & 0.21 & 0.05 & 0.715 & [0.55, 0.87] & 0.736 & 0.513 & 0.541 & 100\% & 83\% \\
Artery\_Tibial & 27 & 542 & 0.11 & 0.05 & 0.744 & [0.64, 0.84] & 0.741 & 0.479 & 0.547 & 100\% & 87\% \\
Cells\_Cultured\_fibroblasts & 11 & 350 & 0.17 & 0.03 & 0.688 & [0.54, 0.82] & 0.639 & 0.584 & 0.597 & 100\% & 82\% \\
Muscle\_Skeletal & 16 & 397 & 0.11 & 0.04 & 0.798 & [0.68, 0.90] & 0.621 & 0.487 & 0.538 & 100\% & 87\% \\
Nerve\_Tibial & 34 & 609 & 0.13 & 0.05 & 0.818 & [0.76, 0.88] & 0.683 & 0.511 & 0.588 & 100\% & 88\% \\
Skin\_Not\_Sun\_Exposed\_Suprapubic & 17 & 456 & 0.09 & 0.04 & 0.828 & [0.73, 0.91] & 0.710 & 0.471 & 0.586 & 100\% & 87\% \\
Skin\_Sun\_Exposed\_Lower\_leg & 26 & 607 & 0.10 & 0.04 & 0.694 & [0.60, 0.78] & 0.700 & 0.544 & 0.574 & 100\% & 91\% \\
Thyroid & 26 & 524 & 0.14 & 0.05 & 0.699 & [0.58, 0.81] & 0.612 & 0.529 & 0.527 & 100\% & 86\% \\
Whole\_Blood & 17 & 305 & 0.18 & 0.05 & 0.776 & [0.66, 0.88] & 0.689 & 0.512 & 0.594 & 100\% & 84\% \\
\bottomrule
\end{tabular}
\end{table}


\section{eQTLP-v2 construction and the full baseline grid}
\label{app:b}

\paragraph{Construction.} We pool the nine tissues into 25,600 unique
variant--gene pairs and drop 111 with conflicting labels across tissues. Each
positive is matched to one negative on log distance to the TSS and on the same
side of the TSS, with a caliper of 0.05 on log distance; this yields 2,905
matched pairs (5,810 rows). Folds hold out whole chromosomes, balanced on
positive count, giving 570--665 test positives per fold. We drop the published
label-derived mask, which marks the positions of the positive eQTLs and is
therefore a label leak.

\paragraph{On matching away distance.} Distance to the TSS is genuinely
predictive of regulatory effect, so matching removes real signal along with the
shortcut, and the matched task is harder than the underlying biology requires.
We match anyway for two reasons. The task's stated aim is variant-effect
prediction from sequence, and a benchmark that a distance baseline wins is not
measuring that. And the published input hands the model the distance directly
through its padding, so the shortcut needs no sequence understanding at all.
Table~\ref{tab:b1} reports the unmatched (\textbf{all}) grid alongside the
matched one so the cost of matching is visible: on the random split with
distance available, the composition baseline reaches 0.70, and the gap between
that and the matched chromosome-holdout floor of 0.599 is what matching removes.

\paragraph{Composition of the matched set.} 5,360 of 5,810 rows (92.3\%) are
single-nucleotide substitutions; the remaining 450 (7.7\%) are short indels or
MNVs, with a median length change of 2\,bp and a maximum of 76\,bp. Indels are
balanced across labels (92.1\% of negatives and 92.4\% of positives are SNVs), so
they cannot carry the label. Nor do they affect the embedding-shift result: the
median relative shift restricted to substitutions is $1.52 \times 10^{-4}$
(HyenaDNA) and $0.674 \times 10^{-4}$ (Caduceus) against $1.52 \times 10^{-4}$
and $0.675 \times 10^{-4}$ over all pairs. Indels shift the representation
slightly more ($1.63$ and $0.691 \times 10^{-4}$), as expected from changing more
bases, but they are too few and too similar to move the pooled statistic.

% generated by analysis/make_appendix_tables.py -- do not edit by hand
\begin{table}[tb]
\centering
\caption{eQTLP-v2 baseline grid. \textbf{all} is every pooled pair; \textbf{matched} is the distance-matched subset the paper uses. \textbf{random} allows locus memorisation; \textbf{gene} and \textbf{chrom} do not.}
\label{tab:b1}
\small
\begin{tabular}{lccccccc}
\toprule
Version & Split & Model & AUROC & fold mean & fold SD & $n$ & $n_+$ \\
\midrule
all & random & distance & 0.790 & 0.790 & 0.008 & 25600 & 3254 \\
all & random & gene\_prior & 0.703 & 0.703 & 0.002 & 25600 & 3254 \\
all & random & composition & 0.744 & 0.744 & 0.006 & 25600 & 3254 \\
all & random & dist+comp & 0.833 & 0.833 & 0.008 & 25600 & 3254 \\
all & random & dist+comp+allele & 0.834 & 0.835 & 0.006 & 25600 & 3254 \\
all & random & dist+mask\_count & 0.810 & 0.811 & 0.007 & 25600 & 3254 \\
all & gene & distance & 0.790 & 0.768 & 0.034 & 25600 & 3254 \\
all & gene & gene\_prior & 0.373 & 0.500 & 0.000 & 25600 & 3254 \\
all & gene & composition & 0.640 & 0.669 & 0.011 & 25600 & 3254 \\
all & gene & dist+comp & 0.788 & 0.779 & 0.031 & 25600 & 3254 \\
all & gene & dist+comp+allele & 0.790 & 0.780 & 0.029 & 25600 & 3254 \\
all & gene & dist+mask\_count & 0.781 & 0.767 & 0.035 & 25600 & 3254 \\
all & chrom & distance & 0.790 & 0.770 & 0.055 & 25600 & 3254 \\
all & chrom & gene\_prior & 0.387 & 0.500 & 0.000 & 25600 & 3254 \\
all & chrom & composition & 0.639 & 0.668 & 0.028 & 25600 & 3254 \\
all & chrom & dist+comp & 0.790 & 0.784 & 0.049 & 25600 & 3254 \\
all & chrom & dist+comp+allele & 0.790 & 0.783 & 0.047 & 25600 & 3254 \\
all & chrom & dist+mask\_count & 0.780 & 0.768 & 0.047 & 25600 & 3254 \\
matched & random & distance & 0.500 & 0.500 & 0.015 & 5810 & 2905 \\
matched & random & gene\_prior & 0.706 & 0.707 & 0.009 & 5810 & 2905 \\
matched & random & composition & 0.665 & 0.665 & 0.013 & 5810 & 2905 \\
matched & random & dist+comp & 0.670 & 0.670 & 0.017 & 5810 & 2905 \\
matched & random & dist+comp+allele & 0.670 & 0.670 & 0.018 & 5810 & 2905 \\
matched & random & dist+mask\_count & 0.529 & 0.529 & 0.014 & 5810 & 2905 \\
matched & gene & distance & 0.500 & 0.488 & 0.039 & 5810 & 2905 \\
matched & gene & gene\_prior & 0.432 & 0.500 & 0.000 & 5810 & 2905 \\
matched & gene & composition & 0.602 & 0.608 & 0.009 & 5810 & 2905 \\
matched & gene & dist+comp & 0.601 & 0.610 & 0.005 & 5810 & 2905 \\
matched & gene & dist+comp+allele & 0.596 & 0.604 & 0.006 & 5810 & 2905 \\
matched & gene & dist+mask\_count & 0.450 & 0.463 & 0.014 & 5810 & 2905 \\
matched & chrom & distance & 0.500 & 0.495 & 0.054 & 5810 & 2905 \\
matched & chrom & gene\_prior & 0.439 & 0.500 & 0.000 & 5810 & 2905 \\
matched & chrom & composition & 0.588 & 0.593 & 0.018 & 5810 & 2905 \\
matched & chrom & dist+comp & 0.582 & 0.590 & 0.029 & 5810 & 2905 \\
matched & chrom & dist+comp+allele & 0.590 & 0.599 & 0.031 & 5810 & 2905 \\
matched & chrom & dist+mask\_count & 0.436 & 0.444 & 0.012 & 5810 & 2905 \\
\bottomrule
\end{tabular}
\end{table}


\section{Per-fold ablation results}
\label{app:c}

Table~\ref{tab:c1} is the per-fold data behind Figure~1a, and
Table~\ref{tab:c2} sweeps the probe's regularisation strength.

The \texttt{ref\_copy} ablation refits nothing: it takes the classifier already
fitted on $[e_{\mathrm{ref}}, e_{\mathrm{alt}}]$ and scores
$[e_{\mathrm{ref}}, e_{\mathrm{ref}}]$, so any difference is attributable to the
alternate allele alone. Across $C \in \{0.001, \ldots, 10\}$, both reading
orientations, linear and MLP heads, and 13 fold partitions, the gap never leaves
the fourth decimal. The paired bootstrap 95\% CI on the difference is
$[-0.0002, 0.0002]$.

\paragraph{Orientation.} HyenaDNA is causal and the published orientation puts
the variant roughly 500\,bp from the end, so in 71--86\% of pairs only a few
hundred positions follow it. Embedding the reverse complement instead, so that
every downstream position can see the variant, changes AUROC by less than 0.01
and leaves the ablation gap unchanged. Bidirectional Caduceus behaves the same
way, which rules out causality as the explanation.

\paragraph{Numerical precision.} Re-embedding 400 pairs in full fp32 reproduces
the bf16 shifts (difference-vector cosine 0.89--0.95 between the two precisions),
and \texttt{diff\_only} stays at chance. The null is not a rounding artifact.

% generated by analysis/make_appendix_tables.py -- do not edit by hand
\begin{table}[tb]
\centering
\caption{Per-fold test AUROC behind Figure 1a, and the \texttt{ref\_copy} ablation on the identical fitted model. The CNN contributes 3 seeds per fold.}
\label{tab:c1}
\small
\begin{tabular}{lcccc}
\toprule
Model & Fold & Seed & \texttt{refalt} & \texttt{ref\_copy} \\
\midrule
CNN & 0 & 0 & 0.4687 & 0.4688 \\
CNN & 0 & 1 & 0.5750 & 0.5752 \\
CNN & 0 & 2 & 0.4915 & 0.4916 \\
CNN & 1 & 0 & 0.5437 & 0.5437 \\
CNN & 1 & 1 & 0.4749 & 0.4756 \\
CNN & 1 & 2 & 0.4861 & 0.4855 \\
CNN & 2 & 0 & 0.5360 & 0.5359 \\
CNN & 2 & 1 & 0.5690 & 0.5694 \\
CNN & 2 & 2 & 0.5025 & 0.5026 \\
CNN & 3 & 0 & 0.5106 & 0.5113 \\
CNN & 3 & 1 & 0.5709 & 0.5713 \\
CNN & 3 & 2 & 0.5255 & 0.5273 \\
CNN & 4 & 0 & 0.4984 & 0.4985 \\
CNN & 4 & 1 & 0.4497 & 0.4492 \\
CNN & 4 & 2 & 0.5596 & 0.5598 \\
Caduceus & 0 & 0 & 0.5492 & 0.5491 \\
Caduceus & 1 & 0 & 0.5874 & 0.5874 \\
Caduceus & 2 & 0 & 0.5350 & 0.5349 \\
Caduceus & 3 & 0 & 0.5772 & 0.5772 \\
Caduceus & 4 & 0 & 0.5686 & 0.5690 \\
Floor (composition) & 0 & 0 & 0.5870 & -- \\
Floor (composition) & 1 & 0 & 0.6144 & -- \\
Floor (composition) & 2 & 0 & 0.5644 & -- \\
Floor (composition) & 3 & 0 & 0.6096 & -- \\
Floor (composition) & 4 & 0 & 0.5905 & -- \\
HyenaDNA & 0 & 0 & 0.5908 & 0.5908 \\
HyenaDNA & 1 & 0 & 0.5225 & 0.5225 \\
HyenaDNA & 2 & 0 & 0.5719 & 0.5717 \\
HyenaDNA & 3 & 0 & 0.5694 & 0.5696 \\
HyenaDNA & 4 & 0 & 0.5250 & 0.5253 \\
\bottomrule
\end{tabular}
\end{table}

\begin{table}[tb]
\centering
\caption{Probe-strength sweep on the fixed chromosome folds. The \texttt{refalt}--\texttt{ref\_copy} gap stays at the fourth decimal at every regularisation strength, and every probe sits below the floor.}
\label{tab:c2}
\small
\begin{tabular}{lcccccc}
\toprule
Model & $C$ & \texttt{refalt} & \texttt{ref\_copy} & \texttt{diff\_only} & floor & $\Delta$ \\
\midrule
HyenaDNA & 0.001 & 0.5704 & 0.5704 & 0.4968 & 0.599 & -0.00003 \\
HyenaDNA & 0.01 & 0.5667 & 0.5667 & 0.4958 & 0.599 & +0.00002 \\
HyenaDNA & 0.1 & 0.5559 & 0.5560 & 0.4953 & 0.599 & -0.00006 \\
HyenaDNA & 1 & 0.5472 & 0.5472 & 0.4955 & 0.599 & -0.00009 \\
HyenaDNA & 10 & 0.5373 & 0.5377 & 0.4955 & 0.599 & -0.00040 \\
Caduceus & 0.001 & 0.5775 & 0.5775 & 0.4972 & 0.599 & -0.00005 \\
Caduceus & 0.01 & 0.5693 & 0.5693 & 0.5031 & 0.599 & -0.00003 \\
Caduceus & 0.1 & 0.5635 & 0.5635 & 0.5068 & 0.599 & -0.00002 \\
Caduceus & 1 & 0.5476 & 0.5477 & 0.5055 & 0.599 & -0.00012 \\
Caduceus & 10 & 0.5341 & 0.5347 & 0.5053 & 0.599 & -0.00059 \\
\bottomrule
\end{tabular}
\end{table}


\section{The CNN collapse}
\label{app:d}

The benchmark's eQTLP head ends in a ReLU applied to the logits. Both logits can
therefore clamp to zero, at which point the prediction is constant and every
weight receives zero gradient: the run is dead and cannot recover. On the
\emph{published} split, 2 of 10 initialisations are dead before the first
optimisation step, and after one epoch 5 of 10 (Whole\_Blood) and 6 of 10
(Nerve\_Tibial) seeds emit a constant score, i.e.\ AUROC exactly 0.5. The
training script sets no random seed and evaluates the final-epoch model rather
than the best checkpoint, so which of these outcomes is reported is a matter of
chance. The CNN is also fed only the reference sequence, so it cannot use the
allele even in principle.

Table~\ref{tab:d1} reports the repaired CNN --- ReLU removed from the logits ---
on eQTLP-v2, across 5 folds and 3 seeds. It sits near chance (0.517 mean) and
below the composition floor, and the allele ablation moves it by $-0.0002$ on
average, with a worst case of 0.0017 across the 15 runs.

% generated by analysis/make_appendix_tables.py -- do not edit by hand
\begin{table}[tb]
\centering
\caption{End-to-end CNN on eQTLP-v2, 5 folds $\times$ 3 seeds. \textbf{alt} is the real pair, \textbf{ref\_copy} replaces the alternate allele with the reference, and \textbf{rand} substitutes a base that never occurred.}
\label{tab:d1}
\small
\begin{tabular}{lccccc}
\toprule
Fold & Seed & alt & \texttt{ref\_copy} & rand & $\Delta$ \\
\midrule
0 & 0 & 0.4687 & 0.4688 & 0.4687 & -0.00006 \\
0 & 1 & 0.5750 & 0.5752 & 0.5753 & -0.00024 \\
0 & 2 & 0.4915 & 0.4916 & 0.4915 & -0.00000 \\
1 & 0 & 0.5437 & 0.5437 & 0.5438 & -0.00003 \\
1 & 1 & 0.4749 & 0.4756 & 0.4751 & -0.00069 \\
1 & 2 & 0.4861 & 0.4855 & 0.4862 & +0.00055 \\
2 & 0 & 0.5360 & 0.5359 & 0.5357 & +0.00008 \\
2 & 1 & 0.5690 & 0.5694 & 0.5693 & -0.00040 \\
2 & 2 & 0.5025 & 0.5026 & 0.5026 & -0.00002 \\
3 & 0 & 0.5106 & 0.5113 & 0.5113 & -0.00073 \\
3 & 1 & 0.5709 & 0.5713 & 0.5710 & -0.00030 \\
3 & 2 & 0.5255 & 0.5273 & 0.5262 & -0.00175 \\
4 & 0 & 0.4984 & 0.4985 & 0.4986 & -0.00008 \\
4 & 1 & 0.4497 & 0.4492 & 0.4499 & +0.00059 \\
4 & 2 & 0.5596 & 0.5598 & 0.5592 & -0.00022 \\
\bottomrule
\end{tabular}
\end{table}


\section{Released-code defects}
\label{app:e}

Found by reading the released repository and running its own loaders on real
records. We did not re-run the published experiments; with these defects the
released code cannot produce the published FM numbers at all, so those numbers
must come from a version other than the one released. We report this as a
reproducibility problem with the artifact.

\begin{enumerate}
\item \textbf{The eQTL loaders return empty arrays.} Both the HyenaDNA and the
Caduceus eQTL datasets pad each input to 450,000 positions and then slice it with
\texttt{[2028500:-2028500]}. Since $2{,}028{,}500 \times 2 > 450{,}000$, the
slice is empty for every record, so the model receives a zero-length sequence.
\item \textbf{The token ids are wrong.} With the slice removed, the loaders take
\texttt{argmax} over a one-hot encoding, which yields ids 0--3. In both
pretrained vocabularies 0--3 are special tokens; the nucleotides A, C, G, T, N
are ids 7--11. The model is fed padding and separator symbols where bases belong.
\item \textbf{The configs do not match the checkpoints.} The shipped configs
specify 2 layers where the released checkpoints have 8 (HyenaDNA-medium-450k) and
16 (Caduceus-Ph-131k), and the pretrained-weights path is left unset, so a run
started from these configs trains a much smaller randomly initialised model.
\item \textbf{The CNN's logit ReLU collapses training}, as in
Appendix~\ref{app:d}.
\end{enumerate}

\section{Clinical set construction and results by stratum}
\label{app:f}

\paragraph{Construction.} We take ClinVar pathogenic and likely-pathogenic SNVs
with a review status of at least one star, annotate them against GENCODE, and
keep only variants that are non-coding, more than 50\,bp from any annotated
splice site, outside every CDS of every transcript, and outside small structural
RNAs. 652 of the survivors map to a gene TSS within 450\,kb. Each is matched to a
benign SNV passing the identical filters, on distance to the TSS and, where
possible, within the same gene and the same location class; 555 of 652 pairs
match on both gene and class.

\paragraph{Why evaluation-only.} After the splice, CDS and structural-RNA
filters, ClinVar leaves 409 regulatory positives at two stars and 1,169 at one
star --- too few, and too concentrated in a handful of genes, to train on
without simply memorising those genes. No model is fitted on clinical labels;
every score is transferred from eQTLP-v2.

\paragraph{Results.} The embedding shift, the transferred eQTLP-v2 probe and
that probe's allele-only component all score 0.489--0.507 over the full set and
0.49--0.52 on the 285 pairs whose pathogenic member has two stars or better. The
only stratum above chance is 5$'$UTR/promoter at 0.53--0.57, and that is a
distance effect: matching there left pathogenic variants closer to the TSS
(median 170\,bp vs 450\,bp), and distance alone scores 0.61 on that stratum.

\paragraph{Gene concentration.} The set is dominated by a few well-studied genes,
so Table~\ref{tab:f1} repeats the evaluation after dropping the three most
frequent genes, after dropping the ten most frequent, and after keeping a single
pair per gene. Every score stays at chance, so the null is not an artifact of
gene composition.

\paragraph{Interpretation.} Pathogenicity and expression effect are different
targets: a variant can shift expression without being pathogenic, or be
pathogenic through a mechanism unrelated to expression. This result therefore
corroborates the eQTL finding on an independent label set rather than showing
that the eQTL conclusion generalises to clinical pathogenicity.

% generated by analysis/make_appendix_tables.py -- do not edit by hand
\begin{table}[tb]
\centering
\caption{Clinical set, sensitivity to gene concentration. Dropping the most frequent genes, or keeping one pair per gene, leaves every score at chance.}
\label{tab:f1}
\small
\begin{tabular}{lcccc}
\toprule
Quantity & all & drop NF1/HBB/ATM & drop top 10 genes & one pair per gene \\
\midrule
pairs & 652 & 564 & 466 & 301 \\
baseline:-dist & 0.51 & 0.51 & 0.504 & 0.5 \\
HyenaDNA:shift\_fwd & 0.503 & 0.501 & 0.496 & 0.495 \\
HyenaDNA:shift\_rev & 0.506 & 0.504 & 0.498 & 0.496 \\
HyenaDNA:eqtl\_probe & 0.503 & 0.506 & 0.497 & 0.495 \\
HyenaDNA:eqtl\_allele & 0.502 & 0.484 & 0.472 & 0.486 \\
Caduceus:shift\_fwd & 0.496 & 0.501 & 0.499 & 0.497 \\
Caduceus:shift\_rev & 0.5 & 0.502 & 0.497 & 0.5 \\
Caduceus:eqtl\_probe & 0.507 & 0.506 & 0.497 & 0.5 \\
Caduceus:eqtl\_allele & 0.489 & 0.483 & 0.482 & 0.488 \\
\bottomrule
\end{tabular}
\end{table}


\end{document}
